% ===========================================================================
%  Chapitre 12 — Dénombrement
%  RAPE, troisième période — outil du calcul des probabilités
% ===========================================================================
\bmtchapitre{Dénombrement}
  {Compter le nombre de résultats d'un tirage : successif avec remise,
   successif sans remise, simultané ; utiliser à bon escient $n^{p}$,
   $\arrang{p}{n}$ et $\combi{p}{n}$.}
  {Traitement des données \textbullet\ environ 6 heures}

Avant de calculer une probabilité, il faut savoir compter. C'est tout l'objet
de ce chapitre, et c'est la difficulté réelle de l'exercice 2 du
baccalauréat : les élèves savent presque tous écrire
$\Prob = \frac{\text{cas favorables}}{\text{cas possibles}}$ ; ce qu'ils
ratent, c'est le décompte de l'un ou de l'autre.

Tout se joue sur trois formules et une seule question : \emph{l'ordre
compte-t-il, et remet-on ce que l'on a tiré ?} Répondre à cette question,
c'est avoir résolu l'exercice ; la formule suit toute seule.

% ---------------------------------------------------------------------------
\begin{revision}
Des acquis de première.

\begin{enumerate}[label=\textbf{\arabic*.}, leftmargin=*, itemsep=0.35em]
  \item Calculer $5\times4\times3$.
  \item Combien de nombres de deux chiffres peut-on écrire avec les
        chiffres $1$, $2$ et $3$, en pouvant les répéter ?
  \item Et si l'on ne peut pas les répéter ?
  \item Calculer $\dfrac{6\times5\times4}{3\times2\times1}$.
  \item Dans une classe de $30$ élèves, de combien de façons choisit-on un
        délégué puis un suppléant ?
  \item Combien y a-t-il de façons de choisir deux élèves parmi trois, sans
        ordre ?
\end{enumerate}
\end{revision}

\begin{automatismes}
Sans calculatrice.

\begin{enumerate}[label=\textbf{\alph*)}, leftmargin=*, itemsep=0.25em]
  \item $3!$
  \item $4!$
  \item $\arrang{2}{5}$
  \item $\combi{2}{5}$
  \item $\combi{1}{7}$
  \item $\combi{7}{7}$
  \item $\combi{0}{9}$
  \item $2^{3}$
\end{enumerate}
\end{automatismes}

\begin{corrige}[automatismes]
a) $6$ \quad b) $24$ \quad c) $20$ \quad d) $10$ \quad e) $7$ \quad
f) $1$ \quad g) $1$ \quad h) $8$.
\end{corrige}

% ===========================================================================
\section{Compter des listes : le principe multiplicatif}

\begin{theoreme}[principe multiplicatif]
Si un choix se fait en $k$ étapes successives, la première offrant $n_{1}$
possibilités, la deuxième $n_{2}$, \dots, la $k$-ième $n_{k}$, alors le
nombre total de choix possibles est
\[
  n_{1}\times n_{2}\times\cdots\times n_{k}.
\]
\end{theoreme}

\begin{propriete}[tirages successifs avec remise]
On tire $p$ objets l'un après l'autre parmi $n$, en \textbf{remettant}
chaque objet avant le suivant. L'ordre compte, et chaque tirage se fait
parmi les $n$ objets :
\[
  N = n^{p}.
\]
\end{propriete}

\begin{exemple}[trois billets tirés avec remise]
Un porte-monnaie contient $12$ billets. On en tire trois, un par un, en
remettant chaque billet avant le suivant. Le nombre de tirages possibles est
$12^{3} = 1\,728$.
\end{exemple}

% ===========================================================================
\section{Arrangements et permutations}

\begin{definition}[factorielle]
Pour un entier $n \geq 1$, on note
$n! = n\times(n-1)\times\cdots\times2\times1$, et l'on pose $0! = 1$.
\end{definition}

\begin{definition}[arrangement]
Un \textbf{arrangement} de $p$ objets parmi $n$ est une liste
\textbf{ordonnée} de $p$ objets \textbf{distincts} choisis parmi $n$. Leur
nombre est
\[
  \arrang{p}{n} = \frac{n!}{(n-p)!}
  = \underbrace{n\times(n-1)\times\cdots\times(n-p+1)}_{p \text{ facteurs}} .
\]
\end{definition}

\begin{remarque}
C'est exactement le décompte d'un \textbf{tirage successif sans remise} : au
premier tirage, $n$ choix ; au deuxième, $n-1$ ; et ainsi de suite. La
formule n'est que le principe multiplicatif écrit une fois pour toutes.
\end{remarque}

\begin{definition}[permutation]
Une \textbf{permutation} de $n$ objets est un arrangement de $n$ objets
parmi $n$, c'est-à-dire un rangement de tous les objets dans un ordre. Leur
nombre est $\arrang{n}{n} = n!$.
\end{definition}

\begin{exemple}[trois billets tirés sans remise]
Dans le même porte-monnaie de $12$ billets, on tire trois billets l'un après
l'autre \emph{sans} remise. Le nombre de tirages possibles est
\[
  \arrang{3}{12} = 12\times11\times10 = 1\,320 .
\]
\end{exemple}

% ===========================================================================
\section{Combinaisons : quand l'ordre ne compte pas}

\begin{definition}[combinaison]
Une \textbf{combinaison} de $p$ objets parmi $n$ est une partie de $p$
objets choisis parmi $n$, \textbf{sans tenir compte de l'ordre}. Leur nombre
est
\[
  \combi{p}{n} = \frac{n!}{p!\,(n-p)!} = \frac{\arrang{p}{n}}{p!} .
\]
\end{definition}

\begin{remarque}
C'est le décompte d'un \textbf{tirage simultané} — « d'un seul coup »,
« en une seule fois », « à la fois » : les expressions employées par les
sujets. On divise par $p!$ parce que les $p!$ ordres possibles d'un même
groupe donnent le même résultat.
\end{remarque}

\begin{propriete}[valeurs utiles]
\[
  \combi{0}{n} = \combi{n}{n} = 1, \qquad
  \combi{1}{n} = n, \qquad
  \combi{p}{n} = \combi{n-p}{n} .
\]
\end{propriete}

\begin{exemple}[deux billets tirés simultanément]
Toujours parmi $12$ billets, on en tire deux en une seule fois :
\[
  \combi{2}{12} = \frac{12\times11}{2} = 66 \text{ tirages possibles}.
\]
\end{exemple}

% ===========================================================================
\section{Reconnaître le bon modèle}

\begin{methode}{deux questions, une formule}
\begin{center}
\renewcommand{\arraystretch}{1.5}
\begin{tabular}{|L{4.3cm}|L{3.1cm}|L{3.5cm}|}
\hline
\textbf{L'énoncé dit} & \textbf{L'ordre} & \textbf{On compte avec}\\
\hline
« successivement, avec remise » & compte & $n^{p}$\\
« successivement, sans remise », « l'un après l'autre » &
compte & $\arrang{p}{n} = \dfrac{n!}{(n-p)!}$\\
« simultanément », « à la fois », « d'un seul coup », « en même temps » &
ne compte pas & $\combi{p}{n} = \dfrac{n!}{p!(n-p)!}$\\
\hline
\end{tabular}
\end{center}
\end{methode}

\begin{visualisation}[le même sac, trois façons de compter]
\begin{center}
\begin{tikzpicture}[scale=0.95]
  % trois colonnes
  \foreach \k/\titre/\res in {0/{avec remise}/{$n^{p}$},
                              1/{sans remise}/{$\arrang{p}{n}$},
                              2/{simultané}/{$\combi{p}{n}$}}{
    \node[font=\footnotesize\bfseries, text=bmtRaphia]
      at (\k*4.5+1.5, 3.05) {\titre};
    \draw[bmtGrisClair, line width=0.8pt, rounded corners=3pt]
      (\k*4.5, -0.55) rectangle (\k*4.5+3.4, 2.7);
    \node[font=\footnotesize] at (\k*4.5+1.7, -0.05) {\res};
  }
  % avec remise : la meme boule revient
  \foreach \i in {0,1,2}{
    \fill[bmtIndigo] (0.7+\i*1.0, 2.05) circle (0.2);
    \node[text=bmtPapier, font=\tiny] at (0.7+\i*1.0, 2.05) {A};}
  \node[font=\scriptsize, text=bmtGris, align=center]
    at (1.7, 1.1) {l'ordre compte,\\ une boule peut revenir};
  % sans remise : trois boules distinctes ordonnees
  \foreach \i/\l in {0/A, 1/B, 2/C}{
    \fill[bmtLaterite] (5.2+\i*1.0, 2.05) circle (0.2);
    \node[text=bmtPapier, font=\tiny] at (5.2+\i*1.0, 2.05) {\l};}
  \node[font=\scriptsize, text=bmtGris, align=center]
    at (6.2, 1.1) {l'ordre compte,\\ pas de répétition};
  % simultane : un groupe
  \draw[bmtVetiver, line width=0.9pt, rounded corners=5pt]
    (9.3, 1.72) rectangle (11.6, 2.38);
  \foreach \i/\l in {0/A, 1/B, 2/C}{
    \fill[bmtVetiver] (9.7+\i*0.75, 2.05) circle (0.2);
    \node[text=bmtPapier, font=\tiny] at (9.7+\i*0.75, 2.05) {\l};}
  \node[font=\scriptsize, text=bmtGris, align=center]
    at (10.6, 1.1) {l'ordre ne compte pas :\\ un seul groupe};
\end{tikzpicture}
\end{center}

Trois tirages de trois objets, trois décomptes différents. À gauche, la
même boule peut sortir plusieurs fois et l'ordre distingue les résultats. Au
milieu, chaque boule ne sort qu'une fois mais l'ordre compte encore : $ABC$
et $CBA$ sont deux tirages. À droite, $\{A,B,C\}$ est un seul résultat,
quelle que soit la façon dont on l'écrit.
\end{visualisation}

\begin{erreurclassique}
Le mot « successivement » seul ne suffit pas : il faut lire la suite de la
phrase. « Successivement \emph{avec} remise » donne $n^{p}$ ;
« successivement \emph{sans} remise » donne $\arrang{p}{n}$. Les deux
formules diffèrent beaucoup : pour $12$ billets et $3$ tirages, $1\,728$
contre $1\,320$.
\end{erreurclassique}

\begin{entrainement}[compter]
\exo[\niveau{1}] Une urne contient $8$ boules. On en tire $3$
simultanément : combien de tirages possibles ?

\exo[\niveau{1}] Même urne, trois tirages successifs avec remise : combien
de tirages possibles ?

\exo[\niveau{2}] Même urne, trois tirages successifs sans remise : combien
de tirages possibles ?

\exo[\niveau{2}] Une trousse contient $12$ stylos. On en tire $3$
successivement sans remise : vérifier qu'il y a $1\,320$ cas possibles.

\exo[\niveau{2}] Dans une classe de $12$ élèves, on choisit $3$ élèves
simultanément : combien de choix ?

\exo[\niveau{3}] Une urne contient $5$ boules noires et $3$ blanches. On en
tire $3$ simultanément. Combien de tirages contiennent exactement deux
noires ?

\exo[\niveau{3}] Avec les huit lettres du mot SCIENCES, on tire quatre
plaquettes simultanément : combien de tirages possibles ?
\end{entrainement}

\begin{corrige}[compter]
\textbf{1.} $\combi{3}{8} = \frac{8\times7\times6}{6} = 56$.
\textbf{2.} $8^{3} = 512$.
\textbf{3.} $\arrang{3}{8} = 8\times7\times6 = 336$.
\textbf{4.} $\arrang{3}{12} = 12\times11\times10 = 1\,320$. \checkmark
\textbf{5.} $\combi{3}{12} = \frac{12\times11\times10}{6} = 220$.
\textbf{6.} On choisit $2$ noires parmi $5$ et $1$ blanche parmi $3$ :
$\combi{2}{5}\times\combi{1}{3} = 10\times3 = 30$.
\textbf{7.} $\combi{4}{8} = 70$.
\end{corrige}

\begin{qcm}
\exo « On tire simultanément $3$ boules parmi $10$ » : le nombre de cas
possibles est :
\textbf{a)} $10^{3}$ \quad \textbf{b)} $\arrang{3}{10}$ \quad
\textbf{c)} $\combi{3}{10}$ \quad \textbf{d)} $3^{10}$ \reponseqcm{c}

\exo $\combi{3}{10}$ vaut :
\textbf{a)} $720$ \quad \textbf{b)} $120$ \quad \textbf{c)} $30$ \quad
\textbf{d)} $1\,000$ \reponseqcm{b}

\exo $\arrang{2}{7}$ vaut :
\textbf{a)} $21$ \quad \textbf{b)} $42$ \quad \textbf{c)} $49$ \quad
\textbf{d)} $14$ \reponseqcm{b}

\exo « Successivement et avec remise, $4$ jetons parmi $6$ » :
\textbf{a)} $6^{4}$ \quad \textbf{b)} $4^{6}$ \quad
\textbf{c)} $\arrang{4}{6}$ \quad \textbf{d)} $\combi{4}{6}$
\reponseqcm{a}

\exo $\combi{8}{10}$ est égal à :
\textbf{a)} $\combi{2}{10}$ \quad \textbf{b)} $\combi{10}{8}$ \quad
\textbf{c)} $\arrang{8}{10}$ \quad \textbf{d)} $80$ \reponseqcm{a}
\end{qcm}

% ===========================================================================
\begin{annales}
Le dénombrement n'est jamais un exercice à lui seul : il ouvre l'exercice de
probabilité. Les extraits ci-dessous en sont les questions de décompte,
telles qu'elles sont posées.

\exoann{Baccalauréat, Madagascar, série A, session 2001 — exercice 2,
question 1.a)}
Un porte-monnaie contient $12$ billets dont $5$ de $500$ Ar, $4$ de
$1\,000$ Ar et $3$ de $2\,000$ Ar. On prend successivement au hasard et sans
remise $3$ billets. Vérifier qu'il y a $1\,320$ cas possibles.

\exoann{Baccalauréat, Madagascar, série A, session 2000 — exercice 1
(extrait)}
Une urne contient $8$ boules indiscernables au toucher dont $4$ blanches et
$4$ noires.
\begin{enumerate}[label=\textbf{\alph*)}, leftmargin=*, itemsep=0.15em]
  \item On tire simultanément $3$ boules : déterminer le nombre de tirages
        possibles.
  \item On effectue $3$ tirages successifs d'une boule, avec remise :
        déterminer le nombre de tirages possibles.
  \item On tire toutes les boules une à une sans remise : déterminer le
        nombre de tirages possibles.
\end{enumerate}

\exoann{Baccalauréat, Madagascar, série A, session 2004 — exercice 2,
question 3.a)}
Une trousse contient douze stylos : $3$ rouges, $4$ verts, $3$ bleus et $2$
noirs. Un élève tire successivement et sans remise trois stylos.
Déterminer le nombre de cas possibles.

\exoann{Baccalauréat, Madagascar, série A, session 2009 — exercice 2,
question 2.a)}
Une urne contient $9$ boules indiscernables au toucher : $2$ vertes
numérotées $1$ et $1$ ; $3$ rouges numérotées $1$, $2$, $3$ ; $4$ blanches
numérotées $1$, $2$, $3$, $4$. On tire successivement au hasard et sans
remise $3$ boules. Démontrer qu'il y a $504$ cas possibles.

\exoann{Baccalauréat, Madagascar, série A, session 2015 — exercice 2,
question 1.a)}
On dispose de huit plaquettes indiscernables au toucher, portant chacune une
lettre du mot « SCIENCES ». On tire simultanément quatre plaquettes :
déterminer le nombre de cas possibles.

\exoann{Baccalauréat, Madagascar, série A, session 2020 — exercice 2,
question 1.a)}
Un sac contient $12$ billets de banque : $6$ de $5\,000$ Ar, $4$ de
$10\,000$ Ar et $2$ de $20\,000$ Ar. On tire au hasard simultanément
$4$ billets. Calculer le nombre de cas possibles.

\exoann{Baccalauréat, Madagascar, série A, session 2022 — exercice 2,
question 2.a)}
Une trousse contient $9$ stylos indiscernables au toucher, dont $5$ bleus et
$4$ rouges. On tire au hasard et successivement avec remise $3$ stylos.
Quel est le nombre de tirages possibles ?
\end{annales}

\begin{corrige}[annales]
\textbf{1.} $\arrang{3}{12} = 12\times11\times10 = 1\,320$. \checkmark

\textbf{2. a)} $\combi{3}{8} = 56$. \quad \textbf{b)} $8^{3} = 512$. \quad
\textbf{c)} Tirer les huit boules une à une sans remise, c'est les
permuter : $8! = 40\,320$.

\textbf{3.} $\arrang{3}{12} = 1\,320$.

\textbf{4.} $\arrang{3}{9} = 9\times8\times7 = 504$. \checkmark

\textbf{5.} $\combi{4}{8} = \frac{8\times7\times6\times5}{24} = 70$.

\textbf{6.} $\combi{4}{12} = \frac{12\times11\times10\times9}{24} = 495$.

\textbf{7.} $9^{3} = 729$.
\end{corrige}

% ===========================================================================
\begin{vraifaux}
\exo $\combi{p}{n}$ est toujours plus petit que $\arrang{p}{n}$.

\exo Un tirage simultané de deux objets et un tirage successif sans remise
de deux objets donnent le même nombre de cas.

\exo $\combi{4}{9} = \combi{5}{9}$.

\exo $0! = 0$.

\exo Dans un tirage avec remise, le nombre de cas possibles peut dépasser le
nombre d'objets de l'urne.
\end{vraifaux}

\begin{corrige}[vrai ou faux]
\textbf{1.} Vrai dès que $p \geq 2$, puisque
$\combi{p}{n} = \frac{\arrang{p}{n}}{p!}$ ; pour $p = 1$ ou $p = 0$, ils
sont égaux.
\textbf{2.} Faux — le second en donne deux fois plus, l'ordre distinguant
les deux écritures d'un même couple.
\textbf{3.} Vrai — $\combi{4}{9} = \combi{9-4}{9} = \combi{5}{9} = 126$.
\textbf{4.} Faux — $0! = 1$, par convention et pour que les formules
restent vraies.
\textbf{5.} Vrai — $8^{3} = 512$ pour seulement huit boules.
\end{corrige}

% ===========================================================================
\begin{situation}[le tirage au sort du tournoi inter-classes]
Un lycée organise un tournoi de football. Douze classes y participent.

\begin{enumerate}[label=\textbf{\arabic*.}, leftmargin=*, itemsep=0.3em]
  \item On tire au sort, en une seule fois, les quatre classes qui joueront
        le premier jour. Combien de tirages différents sont possibles ?
  \item On tire ensuite au sort, l'une après l'autre, les trois classes qui
        occuperont les trois places du podium : première, deuxième,
        troisième. Combien de podiums différents sont possibles ?
  \item Chaque jour du tournoi, on désigne au hasard une classe chargée de
        l'arbitrage, la même pouvant être désignée plusieurs fois. Sur cinq
        jours, combien de désignations différentes sont possibles ?
  \item Expliquer, en une phrase pour chaque question, pourquoi la formule
        n'est pas la même.
\end{enumerate}
\end{situation}

\begin{corrige}[le tirage au sort du tournoi inter-classes]
\textbf{1.} Tirage simultané : $\combi{4}{12} = 495$.

\textbf{2.} Tirage successif sans remise, l'ordre déterminant la place :
$\arrang{3}{12} = 12\times11\times10 = 1\,320$.

\textbf{3.} Tirage successif avec remise : $12^{5} = 248\,832$.

\textbf{4.} Dans la première, l'ordre est sans importance : quatre classes
forment un groupe. Dans la deuxième, l'ordre fait toute la différence entre
le champion et le troisième. Dans la troisième, l'ordre compte \emph{et} la
même classe peut revenir.
\end{corrige}

% ===========================================================================
\begin{jemeteste}
Vingt minutes.

\begin{enumerate}[label=\textbf{\arabic*.}, leftmargin=*, itemsep=0.3em]
  \item Calculer $\arrang{3}{7}$ et $\combi{3}{7}$.
  \item Une urne contient $10$ jetons. Combien de tirages simultanés de $4$
        jetons ?
  \item Même urne : combien de tirages successifs avec remise de $2$
        jetons ?
  \item Une urne contient $6$ boules rouges et $4$ vertes. On tire
        simultanément $3$ boules : combien en contiennent exactement une
        verte ?
  \item De combien de façons peut-on ranger cinq livres sur une étagère ?
\end{enumerate}
\end{jemeteste}

\begin{corrige}[je me teste]
\textbf{1.} $\arrang{3}{7} = 210$ et $\combi{3}{7} = 35$.
\textbf{2.} $\combi{4}{10} = 210$.
\textbf{3.} $10^{2} = 100$.
\textbf{4.} $\combi{1}{4}\times\combi{2}{6} = 4\times15 = 60$.
\textbf{5.} $5! = 120$.
\end{corrige}

% ===========================================================================
\begin{commeaubac}
\bmtsource{Exercice original au format de l'épreuve — série A, options A1 et
A2 \textbullet\ 5 points, partie dénombrement}
Une boîte contient $11$ jetons indiscernables au toucher, numérotés de $1$
à $11$ : $4$ jetons rouges et $7$ jetons blancs. Parmi les rouges, trois
portent un numéro pair ; parmi les blancs, cinq portent un numéro impair.

\begin{enumerate}[label=\textbf{\arabic*.}, leftmargin=*, itemsep=0.25em]
  \item Reproduire et compléter le tableau suivant. \bareme{1 pt}
    \begin{center}\footnotesize
    \renewcommand{\arraystretch}{1.3}
    \begin{tabular}{|l|c|c|c|}
    \hline
    \textbf{Jetons} & \textbf{numéro pair} & \textbf{numéro impair} &
    \textbf{Total}\\
    \hline
    Rouges & $3$ & & $4$\\
    Blancs & & $5$ & $7$\\
    \hline
    \textbf{Total} & & & $11$\\
    \hline
    \end{tabular}
    \end{center}
  \item On tire simultanément $2$ jetons de la boîte. Déterminer le nombre
        de tirages possibles. \bareme{0,5 pt}
  \item Déterminer le nombre de tirages contenant :
    \begin{enumerate}[label=\textbf{\alph*)}, leftmargin=1.4em, itemsep=0.15em]
      \item deux jetons portant chacun un numéro impair ; \bareme{1 pt}
      \item deux jetons rouges portant chacun un numéro pair. \bareme{1 pt}
    \end{enumerate}
  \item On tire maintenant successivement et sans remise $3$ jetons.
        Déterminer le nombre de tirages possibles, puis le nombre de tirages
        commençant par un jeton rouge. \bareme{1,5 pt}
\end{enumerate}
\end{commeaubac}

\begin{corrige}[comme au baccalauréat]
\textbf{1.} Les rouges à numéro impair sont $4-3 = 1$ ; les blancs à numéro
pair sont $7-5 = 2$. Les totaux de colonnes valent donc $3+2 = 5$ pairs et
$1+5 = 6$ impairs, et $5+6 = 11$. \checkmark

\textbf{2.} $\combi{2}{11} = \frac{11\times10}{2} = 55$.

\textbf{3. a)} Il y a $6$ jetons impairs : $\combi{2}{6} = 15$ tirages.
\textbf{b)} Il y a $3$ jetons rouges pairs : $\combi{2}{3} = 3$ tirages.

\textbf{4.} $\arrang{3}{11} = 11\times10\times9 = 990$ tirages possibles.
Ceux qui commencent par un rouge : $4$ choix pour le premier jeton, puis
$\arrang{2}{10} = 90$ pour les deux suivants, soit $4\times90 = 360$.
\end{corrige}

\begin{notepourlaclasse}
Six heures peuvent paraître beaucoup pour trois formules. Elles ne le sont
pas : la difficulté n'est jamais le calcul, toujours la lecture de l'énoncé.
Consacrez la moitié du temps à un seul exercice — lire une phrase, entourer
les mots « simultanément », « successivement », « avec remise », et écrire
la formule avant tout calcul.

Le tableau à double entrée du bloc « comme au baccalauréat » est une
technique que les sujets récents emploient (session 2013 notamment) : il
mérite d'être enseigné pour lui-même, car il transforme un énoncé confus en
un décompte évident.

Pour les classes nombreuses, le tirage au sort du tournoi se prête à une
mise en scène réelle : douze papiers dans un chapeau, et l'on compte
ensemble. Les élèves retiennent la différence entre les trois modèles bien
mieux après l'avoir vue faite.
\end{notepourlaclasse}
